\section{Discussion}\label{sec:discussion}

We set out to make persistent homology of molecular change on an interaction network testable, and then to ask what it measures. Placing the signed paired change on a fixed complex and comparing the two directions of the lower-star filtration gives a statistic that reverses sign exactly when tumour and normal are exchanged (Proposition~\ref{prop:antisym}). That one property gives exact sign-flip tests with the patient as the unit of inference, and it carries over to a network-adjusted score (Proposition~\ref{prop:Z}), to a gene-level participation score and to every directional baseline. The tie rule is needed for this: it orders tied values by a fixed gene index without changing them, so that the two directions are exact mirror images even when measured changes are equal. In our simulations the Level~1 test and the gene family did not show excess rejection under four noise distributions, on development seeds and again on reserved verification seeds.

What the statistic measures is less favourable to a ring-based reading. A natural interpretation of a persistent $\Hone$ class on a network is a ring of changed genes around a gene that behaves differently. In our matched experiment, the pre-declared comparison across data sets favoured no $\Hone$ statistic, and signed Moran's $I$ separated the two arms better. Within matched pairs, the $\Hone$ statistics did respond to exchanging the centre with a neighbour: at the larger effect size the network-adjusted score rose in \pairZTwo\% of pairs. But the exchange also changed local autocorrelation, degree and clustering among the changed genes, and the topology-free statistics B3 and B9 moved in the same direction, by mean amounts similar to that of $A_i$. In both arms the raw asymmetry was strongly negative, which suggests that the lower-change filtration is dominated by cycles that a block of increased genes fills when it enters; we propose this mechanism but have not verified it interval by interval. In a noiseless planted subgraph the two arms differ cleanly, since the module arm's changed set induces a cone with no $\Hone$; on the full noisy network that difference is diluted by many other classes, and we did not measure how much of it survives.

The behaviour of the network-adjusted score under hub change is the second lesson. Relabelling a patient's changes within strata of degree and clustering is a common way to ask whether an arrangement is more than expected from the nodes involved. On the STRING graph, in which the highest-degree genes share \hubHubRatio{} times as many edges as degree-preserving rewiring predicts, the reference did not absorb strong change at those genes. A reference must preserve what the alternative is meant to exclude, here the adjacency among the changed genes, or the test will attribute that adjacency to arrangement; finer degree strata would also be needed. Degree-preserving rewiring of the whole graph is not such a reference, because it does not preserve the subgraph induced by the changed genes and, on this network, removes most triangles.

These findings bear on applications of persistent homology to functions on molecular and other networks \citep{masoomy2021,ramos2025,song2023,abdullahi2025,lee2012}. They do not show that such analyses are wrong. They show that a persistent loop, or a difference in total persistence, should be compared with simpler network statistics before it is interpreted as ring-like structure. The framework here supplies the tools for that comparison: a paired design with an exact test, a matched benchmark, and baselines computed on the same data with the same sign flips. We suggest that studies report, alongside any $\Hone$ finding, the corresponding $\Hzero$, component-size and Moran's $I$ results, and a null model that preserves the adjacency of the genes involved.

Three analyses would sharpen these conclusions. First, the matched experiment should record, for every data set, whether the planted class survives in the full filtration and what lifetime it contributes, so that the part of the paired response due to the planted class can be separated from the rest. Second, whether $\Hone$ adds information beyond the baselines jointly can be tested directly, by comparing a classifier of the two arms built from Moran's $I$, component size and burden with the same classifier after adding $\Hone$, on held-out data sets with both arms of each data set in the same split. Third, extended persistence \citep{cohensteiner2009}, which treats both directions of a function in one diagram and separates classes by type, may isolate hole-like classes more cleanly. Each should be one pre-declared endpoint evaluated on new simulation seeds and reported whatever the outcome; adding several constructions after seeing these results would trade the calibration that is the framework's main strength for a higher chance of a positive finding.

Persistent homology has properties that the baselines lack: it summarises every threshold at once, so its answer does not depend on a cut-off such as the $\tau$ of B1 and B3, and its cycle representatives say where on the network a signal lies. Whether these properties translate into better detection when change sizes differ from patient to patient has not been tested here.

The breast cancer analysis illustrates both lessons. The methylation asymmetry was the strongest topological result. It kept its direction under standardisation, both alternative network thresholds, composition adjustment and the same-plate restriction, but not under rank transformation, and simple counts, component sizes and autocorrelation pointed the same way. A reader who saw only the persistence diagrams might have described it as rings of hypomethylated genes; the baselines are more simply described as more, and more clustered, hypermethylation of the analysed promoters in tumours. The network-adjusted score did not show a departure from the relabelling reference. The RNA result is the reverse case. The raw asymmetry was fragile and changed sign between reasonable analysis choices, while the network-adjusted score was positive in most patients, modest in each, and remained after an exploratory composition adjustment. We read it as evidence that tumour--normal RNA changes are arranged on the interaction network differently from the same changes placed at random positions of similar degree, clustering and expression. We do not read it as evidence of rings: the score is not specific to holes, and connected groups of co-decreased genes, such as vascular, stromal and adipose programmes, would produce it. The composition proxies cover stroma and immune cells but not adipocytes, which are abundant in normal breast, and an independent cohort is needed before the result can be called confirmed. At the gene level, about a third of the RNA participants and about half of the methylation participants were not significant as single genes on equal terms, and about a quarter and a fifth were also missed by neighbourhood smoothing. Such genes are candidates for follow-up because their network neighbourhood changes; whether they matter in their own right is not something these data can show.
